% Generated draft: requires mathematical and source review.
\documentclass[11pt]{article}
\usepackage[margin=25mm]{geometry}
\usepackage{fontspec}
\setmainfont{Latin Modern Roman}
\usepackage{amsmath,amssymb,mathtools,unicode-math}
\setmathfont{Latin Modern Math}
\usepackage{enumitem,needspace,xurl,hyperref}
\setlength{\parindent}{0pt}
\setlength{\parskip}{5pt}
\setlength{\emergencystretch}{3em}
\newcommand{\N}{\mathbb{N}}
\newcommand{\Z}{\mathbb{Z}}
\newcommand{\Q}{\mathbb{Q}}
\newcommand{\R}{\mathbb{R}}
\newcommand{\C}{\mathbb{C}}
\newcommand{\F}{\mathbb{F}}
\newcommand{\A}{\mathbb{A}}
\newcommand{\PP}{\mathbb{P}}
\newcommand{\E}{\mathbb{E}}
\newcommand{\eps}{\varepsilon}
\newcommand{\dd}{\,\mathrm{d}}
\begin{document}
\section*{UnsolvedMath ID: 30001455}
\textbf{Source problem code:} catalogue entry
\subsection{Definitions and mathematical statement}
Let \(\Omega\subset\R^2\) be a bounded strictly convex planar domain with smooth boundary, as stipulated in the supplied record. The record does not specify the precise regularity meant by smoothness or impose an explicit positive-curvature hypothesis.

An ordinary periodic billiard trajectory in \(\Omega\) is a closed polygonal path whose vertices lie on \(\partial\Omega\), whose open edges lie in \(\Omega\), and which obeys specular reflection at each vertex: the tangential component of velocity is preserved and the normal component is reversed. If its successive vertices are \(x_0,\ldots,x_{q-1},x_q=x_0\), its length is
\[
L(\gamma)=\sum_{j=0}^{q-1}\lVert x_{j+1}-x_j\rVert.
\]
The length spectrum refers to the collection of lengths of periodic billiard trajectories. The supplied record does not specify whether repetitions are included, whether only primitive trajectories are intended, or whether multiplicities are retained.

\textbf{Supplied problem, with unresolved notation.} Determine the length spectrum of periodic billiard trajectories in \(\Omega\), including the asymptotic behavior, as \(n\to\infty\), of
\[
T_{m,n}=\sup\{L(\gamma):\gamma\in\Gamma(m,n)\},
\qquad
t_{m,n}=\inf\{L(\gamma):\gamma\in\Gamma(m,n)\}.
\]
These formulas express the supplied notation \(\sup L(\Gamma(m,n))\) and \(\inf L(\Gamma(m,n))\), conditional on specifying the trajectory family \(\Gamma(m,n)\) and its nonemptiness. The source record does not define \(\Gamma(m,n)\), the meanings or admissible ranges of \(m,n\), or whether \(m\) is fixed or varies with \(n\). No convention for empty families is supplied.

\textbf{Limitation of the formulation.} The supplied background explicitly states that the extracted passage is truncated before the actual length-spectrum question. Thus the displayed request preserves the registry wording but does not recover a definite conjecture, asymptotic formula, error estimate, or quantified limiting regime. Those targets cannot be supplied without additional source evidence.
\subsection{Short English statement}
For a bounded strictly convex smooth planar billiard table, the registry asks for the lengths of periodic trajectories and asymptotics of their supremal and infimal lengths in families denoted by \(\Gamma(m,n)\). Those families and the limiting regime are undefined, and the supplied background says the actual source question was truncated. This is therefore an incomplete problem formulation requiring human review.
\subsection{Sources}
\begin{itemize}
\item Supplied local UnsolvedMath registry record, permanent ID \(30001455\), source code \texttt{catalogue entry}, title \emph{Length Spectra of Periodic Convex Billiard Trajectories}, in the collection \emph{Oberwolfach Reports Open Problems}. The record assigns the year 2010 and identifies no individual proposer.
\item Supplied report citation: \emph{Arbeitsgemeinschaft: Mathematical Billards} (2010), Oberwolfach Reports, pp.~955--1015, DOI \texttt{10.4171/owr/2010/17}. \url{https://doi.org/10.4171/owr/2010/17}. This is also the supplied literature source labeled ``original OWR report.'' The report reference is supplied provenance; its contents have not been independently inspected here.
\item Supplied literature source labeled ``Marvizi--Melrose'': \url{https://doi.org/10.4310/jdg/1214437138}. No fuller bibliographic description is supplied.
\item Supplied literature assessment dated \texttt{2026-08-21}: the record labels the problem open, states that the extraction ends before the actual question, and reports that Marvizi--Melrose asymptotics concern trajectories near the boundary without resolving the missing general claim. This is supplied dated curation, not an independently verified literature assessment. The background requests independent verification of time-sensitive status and cited claims.
\item The registry labels the statement \texttt{corrected\_verified}; its accompanying verification note describes only enclosing TeX notation in mathematical delimiters without changing mathematical content. This label does not establish completeness or mathematical correctness.
\end{itemize}
\end{document}
